% CH1 Plain Form (LaTeX)
% Compile with: xelatex ch1-plain-form.tex
\documentclass[12pt,a4paper]{article}

\usepackage[a4paper,margin=1in]{geometry}
\usepackage{fontspec}
\usepackage{polyglossia}
\usepackage{longtable}
\usepackage{array}
\usepackage{enumitem}
\usepackage{setspace}
\usepackage{hyperref}

\setmainlanguage{thai}
\setotherlanguage{english}
\setmainfont{TH Sarabun New}
\newfontfamily\thaifont{TH Sarabun New}

\setlength{\parindent}{0pt}
\setlength{\parskip}{0.45em}
\onehalfspacing

\newcommand{\linefield}[1]{\textbf{#1} \hrulefill}
\newcommand{\boxitem}[1]{\(\square\) #1}
\newcommand{\spacer}{\vspace{0.6em}}

\begin{document}

{\Large\textbf{แบบสำรวจข้อมูลเพื่อจัดทำแผนการพัฒนาองค์กรสุขภาวะข้าราชการพลเรือน}}\\
เวอร์ชันเรียบ (LaTeX) สำหรับพิมพ์/กรอกแบบฟอร์ม

\spacer
\linefield{อีเมลผู้ตอบ}\\[0.6em]
\linefield{ชื่อหน่วยงาน}

\section*{ส่วนที่ 1: ข้อมูลเบื้องต้นของส่วนราชการ}
\linefield{1) ภาพรวมยุทธศาสตร์และทิศทางของส่วนราชการ}\\[2.2em]
\linefield{2) โครงสร้างองค์กรและบทบาทหน้าที่หลัก}\\[2.2em]
\linefield{3) จำนวนข้าราชการรวม}

\subsection*{3.1 จำแนกตามประเภท}
\linefield{ข้าราชการ}\\
\linefield{พนักงานราชการ}\\
\linefield{ลูกจ้าง}\\
\linefield{อื่น ๆ}

\subsection*{3.2 จำแนกตามช่วงอายุ (เฉพาะข้าราชการ)}
\linefield{ไม่เกิน 30 ปี}\\
\linefield{31--40 ปี}\\
\linefield{41--50 ปี}\\
\linefield{51--60 ปี}

\subsection*{3.3 จำแนกตามอายุราชการ}
\linefield{ไม่เกิน 1 ปี}\\
\linefield{1--5 ปี}\\
\linefield{6--10 ปี}\\
\linefield{11--15 ปี}\\
\linefield{16--20 ปี}\\
\linefield{21--25 ปี}\\
\linefield{26--30 ปี}\\
\linefield{31 ปีขึ้นไป}

\subsection*{3.4 จำแนกตามประเภทและระดับตำแหน่ง}
\linefield{ประเภททั่วไป: O1 / O2 / O3 / O4}\\
\linefield{ประเภทวิชาการ: K1 / K2 / K3 / K4 / K5}\\
\linefield{ประเภทอำนวยการ: M1 / M2}\\
\linefield{ประเภทบริหาร: S1 / S2}

\subsection*{4) การลาออกและโอนย้าย (ย้อนหลัง 5 ปี)}
\begin{longtable}{|p{0.30\textwidth}|p{0.13\textwidth}|p{0.13\textwidth}|p{0.13\textwidth}|p{0.13\textwidth}|}
\hline
\textbf{รายการ} & \textbf{2565} & \textbf{2566} & \textbf{2567} & \textbf{2568} \\\hline
จำนวนข้าราชการต้นปี & & & & \\\hline
จำนวนข้าราชการปลายปี & & & & \\\hline
ลาออก/โอนย้าย & & & & \\\hline
ร้อยละ & & & & \\\hline
\end{longtable}

\section*{ส่วนที่ 2: นโยบายและบริบทภายนอกองค์กร}
\linefield{5) นโยบายและยุทธศาสตร์ที่เกี่ยวข้อง}\\[2.2em]
\linefield{6) ข้อมูลบริบทและความท้าทาย (Context \& Challenges)}\\[2.2em]

\section*{ส่วนที่ 3: ข้อมูลสุขภาวะของข้าราชการในปีที่ผ่านมา}
\subsection*{7) สุขภาวะทางกาย}
\boxitem{เฉพาะข้าราชการ} \hspace{2em} \boxitem{บุคลากรทั้งหมด}

\linefield{เบาหวาน}\\
\linefield{ความดันโลหิตสูง}\\
\linefield{โรคหัวใจและหลอดเลือด}\\
\linefield{โรคไต}\\
\linefield{โรคตับ}\\
\linefield{โรคมะเร็ง}\\
\linefield{ภาวะอ้วน/น้ำหนักเกิน}\\
\linefield{อื่น ๆ (จำนวน/ระบุ)}

\subsection*{8) ข้อมูลการลาป่วย}
\boxitem{เฉพาะข้าราชการ} \hspace{2em} \boxitem{บุคลากรทั้งหมด}

\linefield{จำนวนวันลาป่วยรวมต่อปี}\\
\linefield{จำนวนวันลาป่วยเฉลี่ยต่อคนต่อปี}

\subsection*{10) ผลสำรวจสุขภาพจิต}
\boxitem{เฉพาะข้าราชการ} \hspace{2em} \boxitem{บุคลากรทั้งหมด}

\linefield{ภาวะเครียดเรื้อรัง}\\
\linefield{ภาวะวิตกกังวล}\\
\linefield{ปัญหาการนอนหลับ}\\
\linefield{ภาวะหมดไฟ (Burnout)}\\
\linefield{ภาวะซึมเศร้า}

\subsection*{11) ผลสำรวจความผูกพันองค์กร (Employee Engagement)}
\linefield{คะแนนปี 2568}\\
\linefield{คะแนนปี 2567}\\
\linefield{คะแนนปี 2566}\\
\linefield{คะแนนปี 2565}\\
\linefield{คะแนนปี 2564}\\
\linefield{ประเด็นที่ได้คะแนนต่ำ}

\subsection*{12) ผลสำรวจสุขภาวะด้วยเครื่องมืออื่น ๆ}
\linefield{รายละเอียด}\\[2.2em]

\section*{ส่วนที่ 4: ระบบการบริหารและสภาพแวดล้อมที่เอื้อต่อสุขภาวะ}
\subsection*{13) สถานะระบบการบริหารบุคลากร (เฉพาะข้าราชการ)}
สำหรับแต่ละหัวข้อ โปรดระบุ: \boxitem{มีตามแผน} \boxitem{มีไม่ครบตามแผน} \boxitem{ไม่มี}

\linefield{ระบบพี่เลี้ยง (Mentoring)}\\
\linefield{ระบบหมุนเวียนงาน (Job Rotation)}\\
\linefield{แผนพัฒนารายบุคคล (IDP)}\\
\linefield{เส้นทางความก้าวหน้า (Career Path)}

\subsection*{14) ชั่วโมงการอบรมเฉลี่ยต่อคนต่อปี}
\linefield{ปี 2568}\\
\linefield{ปี 2567}\\
\linefield{ปี 2566}\\
\linefield{ปี 2565}\\
\linefield{ปี 2564}

\subsection*{16) การจัดสภาพแวดล้อมตามหลักการยศาสตร์ (Ergonomics)}
\boxitem{ยังไม่มี} \hspace{1.5em}
\boxitem{มีแผนแต่ยังไม่ดำเนินการ} \hspace{1.5em}
\boxitem{อยู่ระหว่างดำเนินการ} \hspace{1.5em}
\boxitem{ดำเนินการแล้ว}

\linefield{รายละเอียด/ผลลัพธ์}\\[2.2em]

\section*{ส่วนที่ 5: ทิศทาง เป้าหมาย และข้อเสนอแนะ}
\subsection*{17) จุดเน้นการพัฒนา (เลือกไม่เกิน 3 ข้อ พร้อมจัดอันดับ)}
\linefield{อันดับ 1}\\
\linefield{อันดับ 2}\\
\linefield{อันดับ 3}

\subsection*{18) ข้อเสนอแนะเกี่ยวกับ Intervention Packages เพิ่มเติม}
\linefield{รายละเอียด}\\[2.2em]

\subsection*{19) ไฟล์ HRD PLAN ปี 2567--2570 พร้อมผลการปฏิบัติงานตามแผนในปีล่าสุด}
\linefield{สรุปผลการปฏิบัติงาน}\\[2.2em]
\linefield{หมายเหตุ/รายการเอกสารแนบ}

\vfill
{\small จัดทำโดยสถาบันบัณฑิตพัฒนบริหารศาสตร์}

\end{document}
